\scanpage{96}
\mspage{95}
l'extension $\mathcal E_{C'_{\unsure{\xi}}}$ de $\mu$ par $\pi$, et l'extension
$\mathcal E$, avec une condition de passage d'un
tel isomorphisme à un autre à préciser,
pour $\unsure{\xi_1,\xi_2}\in\operatorname{Ob}C$ dans une même composante
connexe de $C$.
\margin{\uline{NB} On récupère $C$ à partir de $C'$
avec la $\mathcal E(\pi,\mu)$ structure,
\deleted{\unsure{le quotient défini}}
\unsure{on pose}\quad $C=C'/\mathcal E(\pi,\mu)$.}

Pour préciser ce point, je vais m'inspirer
de l'exemple précédent avec $C',C$
décrits en termes de $P,T$ et des
\unsure{homomorphismes}
\[
E\longrightarrow P
\qquad
\mu\mathbin{\text{---}}U,
\]
en notant que l'on a alors un diagramme
commutatif
\diagram{\begin{tikzcd}[ampersand replacement=\&]
E \arrow[r,no head] \arrow[d] \& P \arrow[d,"q"] \\
E/\mu=E_0 \arrow[r] \& P_0=\unsure{P/\mu}
\end{tikzcd}}\quad(18)
\margin{On suppose que $\mu\hookrightarrow U$
et $\mu$ distingué dans $U$}
Notons que \deleted{\unsure{pour le}} $P_0$ est un $U_0$-torseur sur $T$,
avec
\[
(19)\qquad U_0=U/\mu
\]
et par suite on a $\widetilde U_0=\widetilde U$ et un iso can
\[
(20)\qquad \unsure{\mathcal E}\simeq\pi_1(U_0,1_{U_0}).
\]
Par suite, \deleted{\unsure{par deux}} $\Pi_1(P_0)$ est muni d'une
